\clearpage

\section{绪论}

\subsection{研究背景与意义}

\par 随着大数据与人工智能技术的快速发展，数据已成为推动社会与经济发展的重要资源。
其中，图（Graph）作为一种能够刻画实体及其相互关系的数据结构，被广泛应用于社交网络分析、生物信息学、推荐系统、
知识图谱以及金融风控等多个领域\cite{hamilton_representation_2018}。在这些应用场景中，节点
通常表示实体（如用户、蛋白质或商品），而边则表示实体之间的关系（如好友关系、相互作用或交易关系）。
通过对图数据进行分析，可以挖掘出潜在的结构模式和规律，为决策提供重要支持\cite{jin_dabigdata_privacy_2020,chen_bigdata_security_2017,zhao_dp_survey_2023}。

\par 然而，图数据中的边连接关系往往直接反映实体之间的私密互动。例如在社交关系图中观察到节点 A 与 B 相连，
即意味着两者之间存在某种社交关联；在邮件往来图中观察到一条边，则意味着两个账号之间存在过通信记录。
已有研究表明，攻击者一旦获取原始图，可通过结构挖掘、属性匹配以及辅助信息融合等手段对节点进行重识别，
进而推断特定用户的社交圈、政治倾向、健康状况等敏感信息。传统的隐私保护方法主要包括数据匿名化、
$k$-匿名以及扰动技术等\cite{1018213520.nh,sweeney_k-anonymity_2002,narayanan_-anonymizing_2009,JSJX202502009}，这些方法通常通过删除或模糊节点标识信息以降低隐私泄露风险。然而大量研究表明，
此类方法在面对背景知识攻击时较为脆弱，攻击者可结合外部信息对匿名数据进行再识别，难以提供严格的隐私保障\cite{1018891994.nh}。

\par 差分隐私（Differential Privacy, DP）由 Dwork 等人于 2006 年提出，是当前公认的具备严格数学定义的隐私保护标准\cite{hutchison_differential_2006,dwork_algorithmic_2014,mcsherry_mechanism_2007}。
其核心思想是：通过向算法的输出注入精心设计的随机噪声，使得任何单条记录的存在与否都不会显著改变输出的概率分布，
从而限制攻击者通过观察输出反推个体信息的能力。在图数据场景下，研究者将差分隐私扩展为节点差分隐私（node-DP）
和边差分隐私（edge-DP）两种典型变体：前者以单个节点的存在与否为邻接条件，后者以单条边的存在与否为邻接条件。
其中，边差分隐私的隐私粒度与实际场景中的连接关系直接对应，且更易于实现可用性较好的算法，
因此成为当前研究和应用中较为主流的模型。

\par 围绕差分隐私下的图数据发布，学术界已经提出了一系列方法。早期的方法如 TmF 直接对邻接矩阵的每个单元加噪，
方法简单但噪声规模随节点数二次增长，可用性较为有限\cite{JSJX202204002,JSJX202502009,1025785867.nh}；DER 通过节点重排和四叉树估计稠密区域，但树深较大、
空间复杂度较高，难以扩展到大图\cite{chen_correlated_2014}；PrivHRG 利用层次随机图编码原图结构，能够保留路径信息但训练开销较大\cite{xiao_differentially_2014}；
LDPGen 借助 $k$-means 聚类降维，但聚类质量在噪声扰动下难以保证\cite{qin_generating_2017}。袁权等人于 2023 年提出的 PrivGraph 方法利用
图的社区结构作为信息抽取与噪声注入的中间粒度，在隐私预算与可用性之间取得了较好的平衡，
并在多个公开数据集和图度量指标上较上述基线方法有所改善，是当前差分隐私图发布方向中具有代表性的工作之一\cite{yuan_privgraph_2023}。

\par 尽管 PrivGraph 在已有方法中表现较为良好，但其图重建阶段仍存在一定的优化空间。分析表明，
原算法在生成跨社区边时采用均匀随机采样，且未对采样过程进行碰撞检测，使得实际生成的边数可能因重复采样而低于预算。
此外，原算法基于 Chung-Lu 模型重建社区内部边，该模型主要保留了节点度数的偏好性，
对真实图中常见的局部传递性结构（即三角形结构）刻画相对有限，因而合成图的聚类系数往往低于原图；同时，
由于跨社区边在生成过程中仅受边数约束、缺乏其他结构性引导，合成图在模块度等社区结构指标上与原图也存在一定差距。
值得注意的是，图重建阶段不消耗任何额外的隐私预算（属于差分隐私的后处理过程），
因此对该阶段进行算法层面的优化不会改变 PrivGraph 的整体隐私性，却能在一定程度上提升合成图的可用性。
本文正是从这一角度出发，对 PrivGraph 的图重建阶段进行改进。

\subsection{国内外研究现状}

\subsubsection{差分隐私下的图统计发布}

\par 差分隐私下的图统计任务是该领域早期研究的重点。Hay 等人提出了基于约束推断的度分布发布算法，
通过对噪声后的度序列进行单调约束修正，提高了度分布的精度\cite{hay_accurate_2009}；Karwa 等人研究了三角形和子图计数问题，
提出利用平滑敏感度（smooth sensitivity）降低噪声规模的方法\cite{karwa_private_2014}；Raskhodnikova 和 
Smith 给出了 Lipschitz 扩展技术，用于将节点级敏感度有界化，使节点差分隐私下的图统计成为可能\cite{raskhodnikova_lipschitz_2016}；
Sun 等人则在去中心化差分隐私下研究了子图计数问题。上述方法对单一统计量进行了精细化设计，
但其输出仅限于特定的统计量，难以直接支持任意的下游分析任务\cite{sun_analyzing_2019}。

\subsubsection{差分隐私下的图发布}

\par 与统计任务不同，图发布的目标是输出一张完整的合成图，
使数据使用者可在该图上自由开展任意分析任务。
Nguyen 等人\cite{nguyen_differentially_2015} 提出了 TmF（Top-m Filter）方法，
对邻接矩阵中的每个元素加拉普拉斯噪声后选取前 $m$ 个最大值作为边，
方法简单但噪声开销随节点数二次增长，在小预算下保留真实图结构的能力较为有限。
Chen 等人\cite{chen_correlated_2014} 提出的 DER（Density-based Exploration and Reconstruction）
通过节点重排集中边到稠密区域，再用噪声四叉树估计密度，
可在中等规模的图上取得相对较好的效果，
但其时间和空间复杂度均为节点数的二次量级，对大图的适用性较差。
Xiao 等人\cite{xiao_differentially_2014} 提出的 PrivHRG 利用层次随机图（HRG）编码原图，
通过 MCMC 采样获得高似然的 HRG 后再加噪发布；
该方法保留了路径信息，但 MCMC 收敛较慢，运行时间显著高于其他方法。
Qin 等人\cite{qin_generating_2017} 在本地差分隐私（LDP）模型下提出了 LDPGen，
使用 $k$-means 聚类对节点降维，可扩展到大图，但聚类质量在噪声扰动下不易保证。

\par 袁权等人于 USENIX Security 2023 上提出的 PrivGraph 是该方向中较为有代表性的工作。
该方法将社区作为信息抽取的中间粒度，在抽取社区内度序列和跨社区边数后再注入噪声并重建图，
从而在大颗粒度的统计稳健性与小颗粒度的结构精细度之间取得了相对平衡。已有实验表明，
PrivGraph 在多个公开数据集和度量指标上较 TmF、DER、PrivHRG 和 LDPGen 等方法均有所改善，
是本文工作的直接基础。

\subsubsection{图数据合成的其他相关研究}

\par 除上述差分隐私下的图发布方法外，图数据合成在其他方向上也有若干相关工作
基于深度生成模型的方法（如 GraphGAN\cite{wang_graphgan_2018}、
GraphVAE\cite{kurkova_graphvae_2018}、NetGAN\cite{bojchevski_netgan_2018} 等）
能够学习图的分布并生成新的样本，但通常缺乏严格的隐私保证；
将差分隐私与生成对抗网络相结合的 DP-GAN\cite{xie_differentially_2018} 等方法
在表格数据上取得了一定效果，但在图数据上的应用仍处于探索阶段，
性能尚不及面向图设计的专用方法。
Zhang 等人提出的 PrivDPR\cite{zhang_privdpr_2025} 方法将 PageRank 机制与深度模型相结合，
在多层结构设计与权重归一化技术的基础上引入差分隐私噪声，
实现了节点级差分隐私下的图生成。在 LDP 设置下，
Wei 等人的 AsgLDP\cite{wei_asgldp_2020}、
Du 等人的 LDPTrace\cite{du_ldptrace_2023} 等方法
将 LDP 扩展到带属性的图和轨迹数据，
进一步丰富了差分隐私图发布的研究图景。

\subsection{本文主要工作}

\par 本文以 PrivGraph 为基础，针对其图重建阶段进行优化研究，主要工作可归纳为以下两点：

\begin{enumerate}
    \item 提出了针对跨社区边处理的图重建优化方法。本文首先对 PrivGraph 图重建阶段的局限性进行了分析：
    原算法在生成跨社区边时按 $\hat{v}$给定的边数预算进行均匀随机采样，
    实际生成的跨社区边数往往与预算接近甚至更多，这看似"还原"了原图的边数，
    实则在合成图中引入了过量的跨社区连接。由于跨社区边在 Louvain 等模块度优化算法中
    具有"桥接"作用，过多的跨社区边会显著扰动后续的社区划分结果，
    使得对合成图重新执行社区检测得到的划分与原图差异较大，从而拉低 NMI 等社区结构保持指标。
    考虑到 PrivGraph 中社区划分本身就受到差分隐私噪声扰动而存在偏差，
    在此基础上若再忠实生成所有跨社区边，会进一步放大这一偏差。
    
    基于这一观察，本文对图重建阶段做出如下改动：
    在生成跨社区边时不再按 $\hat{v}$ 全量采样，而是仅以 10\% 的比例
    保留少量跨社区连接以维持图的整体连通性；为补偿因此引起的总边数下降并保持节点度分布的合理性，
    在社区内部按 5\% 的比例追加少量边。该机制以"牺牲少量跨社区边的真实性、
    换取社区结构的稳定性"为核心，使合成图在重新进行社区检测时能更好地保持
    与原图一致的社区划分。该改动仅作用于已发布的噪声中间结果
    （$\hat{d}$、$\hat{v}$ 等），不再访问原始数据，
    根据差分隐私的后处理性质，不消耗额外的隐私预算。
    
    \item 引入了一个边交换后处理模块作为补充。
    该模块作用于已合成的图上，分两步进行调整：
    
    第一步是\textbf{社区内部边交换}。在每个社区中识别``孤立边''
    （两端节点没有公共邻居的边），将其替换为两端拥有较多公共邻居的边。
    其原理在于：聚类系数本质上度量``一个节点的两个邻居之间是否相连''
    （即是否构成三角形），孤立边对所在邻域的聚类系数无贡献，
    而将边迁移到``两端共享邻居''的位置后，新边更易参与三角形结构，
    从而提升合成图的局部聚类系数。该交换对节点度分布产生轻微扰动，
    实验显示这种扰动对整体度分布的影响可以忽略。
    
    第二步是\textbf{跨社区边的选择性清理}。对每条跨社区边 $(u,v)$ 计算
    其在零模型下的期望连接强度 $\frac{k_u k_v}{2m}$，
    保留期望强度大于 $1$ 的边（这些边在度分布意义下"本来就该存在"，
    属于自然连接），从期望强度较低的边中随机删除一部分。
    其原理在于：低期望强度的跨社区边更可能是差分隐私噪声引入的``虚假桥梁'',
    在合成图上重新执行 Louvain 等模块度优化算法时，
    这类边会诱导本应分离的社区被合并，扰乱社区划分结果；
    清理这部分边可以让合成图的拓扑社区结构更清晰，
    从而提升后续社区检测的稳定性。
    
    该模块同样仅访问合成图与已发布的社区划分结果，
    符合差分隐私的后处理性质，不消耗额外的隐私预算。
    实验结果显示，该模块在部分指标上能够提供一定程度的改善，
    与第一项工作形成互补。
\end{enumerate}

\par 本文在多个数据集如 Chameleon 、CA-HepPh 、Facebook上对所提方法进行了系统
的实验评估，对比了 7 个常用图度量指标在 7 个隐私预算下的表现。实验结果表明，
本文方法在 NMI、模块度相对误差等指标上相对原始 PrivGraph 取得了一定的改善，
在低隐私预算下改善幅度相对更大；在度分布 KL 散度、特征向量中心性等指标上
与原算法保持一致或略有改善，未出现明显的性能回退。消融实验进一步表明，
所观察到的性能提升主要来自跨社区边处理的优化，边交换后处理对部分指标提供了
进一步的改善。

\subsection{论文组织结构}

\par 本文的组织结构如下：

\par 第一章为绪论，介绍研究背景与意义、国内外研究现状、本文主要工作以及论文组织结构。

\par 第二章为相关理论基础，介绍差分隐私的定义和常用机制、图与社区检测的基本概念，并对 PrivGraph 算法的三阶段流程进行回顾。

\par 第三章对 PrivGraph 图重建阶段进行系统的局限性分析，从理论与实现两个层面剖析其存在的问题，为本文的优化设计提供动机。

\par 第四章详细介绍本文提出的两个优化模块：图重建优化与边交换后处理，给出每个模块的设计动机、核心步骤、伪代码以及隐私性证明。

\par 第五章为实验与结果分析，介绍实验设置与评估指标，并通过 7 个度量指标在 7 个隐私预算下的对比展示优化方法的效果。

\par 第六章为总结与展望，对本文的主要贡献进行归纳，分析仍然存在的不足，并对未来的研究方向进行展望。